\documentclass[border=2pt]{standalone}
\usepackage{amsmath, amsthm, amsfonts}
\usepackage{tikz-cd, kotex}
\usepackage{quiver}
\usepackage{Operators}
\usepackage{palette}
\usepackage{diagram-style}
\begin{document}

%%FIG a map breaking effectivity when a-b is at least 2 (Research_Note_2026_10_29-1.svg)
\begin{tikzpicture}[line cap=round, line join=round]
  %% domain: Z covering E (a-b) times, with an a-fold cover C_1 of F_p attached at an unramified point over p
  \draw[accent1, line width=1.2pt] (-1.3,.5)--(1.3,.5);
  \node[accent1, above] at (0,.5) {$Z$};
  \node[accent1, right, font=\scriptsize] at (1.3,.5) {$(a-b){:}1$};
  \draw[line width=.8pt] (.3,.5)--(.3,-.9);
  \fill (.3,.5) circle[radius=1.3pt];
  \node[right, font=\scriptsize] at (.3,.25) {$1$};
  \node[right, font=\scriptsize] at (.3,-.3) {$a{:}1$};
  \node[below] at (.3,-.9) {$C_1$};
  %% the map
  \draw[thick, -{stealth}] (2.5,0)--(3.9,0);
  \node[above] at (3.2,0) {$f$};
  %% image: E and the fiber F_p
  \begin{scope}[shift={(5.6,0)}]
    \draw[accent1, line width=1.2pt] (-1.3,0)--(1.3,0);
    \node[accent1, right] at (1.3,0) {$E$};
    \node[accent1, right, font=\scriptsize] at (1.3,-.35) {$\times(a-b)$};
    \draw[line width=.8pt] (.3,-1)--(.3,1);
    \node[above] at (.3,1) {$F_p$};
    \node[right, font=\scriptsize] at (.3,.55) {$\times a$};
    \fill[accent1] (.3,0) circle[radius=2pt];
    \node[accent1, below left] at (.3,-.05) {$p$};
  \end{scope}
\end{tikzpicture}

%%FIG a stable map of class beta_m: a section component with components covering fibers and a contracted component (Research_Note_2026_10_29-2.svg)
\begin{tikzpicture}[line cap=round, line join=round]
  %% domain: a section component S of class beta_c, components covering fibers, a contracted component
  \draw[line width=1.2pt] (-1.7,.4)--(1.7,.4);
  \node[above] at (-1.3,.4) {$S$};
  \node[right, font=\scriptsize] at (1.7,.4) {$\beta_c$};
  \foreach \x/\d in {-.7/{d_1}, .3/{d_2}}{
    \draw[line width=.8pt] (\x,.4)--(\x,-1);
    \fill (\x,.4) circle[radius=1.3pt];
    \node[right, font=\scriptsize] at (\x,-.45) {$\d{:}1$};
  }
  \draw[black!50, line width=.8pt] (1.2,.4)--(1.2,1.5);
  \fill (1.2,.4) circle[radius=1.3pt];
  \foreach \y/\l in {.9/{x_1}, 1.3/{x_2}}{
    \fill[accent1] (1.2,\y) circle[radius=1.6pt];
    \node[accent1, right, font=\scriptsize] at (1.2,\y) {$\l$};
  }
  %% the map
  \draw[thick, -{stealth}] (2.6,0)--(4,0);
  \node[above] at (3.3,0) {$f$};
  %% image: fibers of rho (vertical), E (horizontal), f(S) a section crossing E c times
  \begin{scope}[shift={(6.6,0)}]
    \foreach \x in {-2,-1.6,-1.2,-.8,-.4,0,.4,.8,1.2,1.6,2}{\draw[black!25, line width=.5pt] (\x,-1.1)--(\x,1.7);}
    \draw[accent1, line width=1.2pt] (-2.2,-.3)--(2.2,-.3);
    \node[accent1, right] at (2.2,-.3) {$E$};
    \draw[line width=1.2pt] plot[domain=-2.1:2.1, samples=80] (\x,{.35+.9*sin(deg(1.4*\x+.3))});
    \node[right] at (2.1,{.35+.9*sin(deg(3.24))}) {$f(S)$};
    \foreach \x in {-1.88,-.79}{\fill[accent1] (\x,-.3) circle[radius=1.8pt];}
    \foreach \x/\d in {-.4/{d_1}, .8/{d_2}}{
      \draw[line width=.9pt] (\x,-1.1)--(\x,1.7);
      \fill (\x,{.35+.9*sin(deg(1.4*\x+.3))}) circle[radius=1.6pt];
      \node[above, font=\scriptsize] at (\x,1.7) {$\times\d$};
    }
    \fill[black!50] (1.6,{.35+.9*sin(deg(2.54))}) circle[radius=2pt];
    \node at (0,-1.6) {$c+d_1+d_2=m$};
  \end{scope}
\end{tikzpicture}

%%FIG the two shapes of combs removed at step j for beta_m: center covering E once, center contracted to p (Research_Note_2026_10_29-3.svg)
\begin{tikzpicture}[line cap=round, line join=round]
  %% (a) center Z covers E once, all branches cover fibers; contact orders sum to j+1
  \begin{scope}
    \draw[accent1, line width=1.2pt] (-1.7,.4)--(1.7,.4);
    \node[accent1, above] at (.2,.4) {$Z$};
    \node[accent1, right, font=\scriptsize] at (1.7,.4) {$1{:}1$};
    \draw[accent1, fill=white] (-1.3,.4) circle[radius=2pt];
    \node[accent1, below] at (-1.3,.35) {$x_n$};
    \foreach \x/\c/\d/\l in {-.4/{m_1}/{d_1}/{C_1}, .8/{m_2}/{d_2}/{C_2}}{
      \draw[line width=.8pt] (\x,.4)--(\x,-1);
      \fill (\x,.4) circle[radius=1.3pt];
      \node[right, font=\scriptsize] at (\x,.15) {$\c$};
      \node[right, font=\scriptsize] at (\x,-.4) {$\d{:}1$};
      \node[below] at (\x,-1) {$\l$};
    }
    \draw[thick, -{stealth}] (2.6,0)--(4,0);
    \node[above] at (3.3,0) {$f$};
    \begin{scope}[shift={(6.6,0)}]
      \draw[accent1, line width=1.2pt] (-2.2,-.2)--(2.2,-.2);
      \node[accent1, right] at (2.2,-.2) {$E$};
      \foreach \x/\d/\p in {-.4/{d_1}/{p_1}, .8/{d_2}/{p_2}}{
        \draw[line width=.9pt] (\x,-1.1)--(\x,1.3);
        \node[above, font=\scriptsize] at (\x,1.3) {$\times\d$};
        \fill[accent1] (\x,-.2) circle[radius=1.8pt];
        \node[accent1, below right] at (\x,-.22) {$\p$};
      }
      \draw[accent1, fill=white] (-1.5,-.2) circle[radius=2pt];
      \node[accent1, above] at (-1.5,-.15) {$f(x_n)$};
    \end{scope}
    \node at (3.3,-1.75) {(a)\quad$m_1+m_2=j+1$};
  \end{scope}
  %% (b) center Z contracted to p, a section branch S of class beta_c and a branch covering F_p; contact orders sum to j
  \begin{scope}[shift={(0,-4.3)}]
    \draw[accent1, line width=1.2pt] (-1.7,.4)--(1.7,.4);
    \node[accent1, above] at (.2,.4) {$Z\mapsto p$};
    \draw[accent1, fill=white] (-1.3,.4) circle[radius=2pt];
    \node[accent1, below] at (-1.3,.35) {$x_n$};
    \draw[line width=1.2pt] (-.4,.4)--(-.4,-1);
    \fill (-.4,.4) circle[radius=1.3pt];
    \node[right, font=\scriptsize] at (-.4,.15) {$m_0$};
    \node[below] at (-.4,-1) {$S$};
    \node[right, font=\scriptsize] at (-.4,-.4) {$\beta_c$};
    \draw[line width=.8pt] (.8,.4)--(.8,-1);
    \fill (.8,.4) circle[radius=1.3pt];
    \node[right, font=\scriptsize] at (.8,.15) {$m_1$};
    \node[right, font=\scriptsize] at (.8,-.4) {$d_1{:}1$};
    \node[below] at (.8,-1) {$C_1$};
    \draw[thick, -{stealth}] (2.6,0)--(4,0);
    \node[above] at (3.3,0) {$f$};
    \begin{scope}[shift={(6.6,0)}]
      \foreach \x in {-2,-1.6,-1.2,-.8,-.4,0,.4,.8,1.2,1.6,2}{\draw[black!25, line width=.5pt] (\x,-1.1)--(\x,1.3);}
      \draw[accent1, line width=1.2pt] (-2.2,-.2)--(2.2,-.2);
      \node[accent1, right] at (2.2,-.2) {$E$};
      \draw[line width=1.2pt] plot[domain=-2.1:2.1, samples=80] (\x,{-.2+.35*\x+.05*\x*\x*\x});
      \node[right] at (2.1,1) {$f(S)$};
      \draw[line width=.9pt] (0,-1.1)--(0,1.3);
      \node[above, font=\scriptsize] at (0,1.3) {$\times d_1$};
      \fill[accent1] (0,-.2) circle[radius=2pt];
      \node[accent1, below right] at (0,-.25) {$p=f(x_n)$};
    \end{scope}
    \node at (3.3,-1.75) {(b)\quad$m_0+m_1=j$};
  \end{scope}
\end{tikzpicture}

%%FIG maps with x_n on E for beta_1 at the first step (Research_Note_2026_10_29-4.svg)
\begin{tikzpicture}[line cap=round, line join=round]
  %% (a) E with two distinct fibers: stays
  \begin{scope}
    \draw[accent1, line width=1.2pt] (-1.3,3.4)--(1.3,3.4);
    \node[accent1, above] at (0,3.4) {$Z\to E$};
    \foreach \x/\l in {-.1/{F_p}, .7/{F_{p'}}}{
      \draw[line width=.8pt] (\x,3.4)--(\x,2.4);
      \fill (\x,3.4) circle[radius=1.3pt];
      \node[right, font=\scriptsize] at (\x,3.15) {$1$};
      \node[below] at (\x,2.4) {$\l$};
    }
    \draw[accent1, fill=white] (-.95,3.4) circle[radius=2pt];
    \node[accent1, below] at (-.95,3.35) {$x_n$};
    \draw[->, black!50] (1.6,2.8)--(1.6,1.3) node[midway, right, text=black] {$f$};
    \draw[accent1, line width=1.2pt] (-1.4,0)--(1.4,0);
    \node[accent1, right] at (1.4,0) {$E$};
    \foreach \x/\l in {-.1/{F_p}, .7/{F_{p'}}}{
      \draw[line width=.8pt] (\x,-1)--(\x,1);
      \node[above] at (\x,1) {$\l$};
    }
    \draw[accent1, fill=white] (-.95,0) circle[radius=2pt];
    \node at (0,-1.6) {(a)};
  \end{scope}
  %% (b) E with a double cover of one fiber: comb
  \begin{scope}[shift={(4.2,0)}]
    \draw[accent1, line width=1.2pt] (-1.3,3.4)--(1.3,3.4);
    \node[accent1, above] at (0,3.4) {$Z\to E$};
    \draw[line width=.8pt] (.3,3.4)--(.3,2.4);
    \fill (.3,3.4) circle[radius=1.3pt];
    \node[right, font=\scriptsize] at (.3,3.15) {$1$};
    \node[right, font=\scriptsize] at (.3,2.8) {$2{:}1$};
    \node[below] at (.3,2.4) {$F_p$};
    \draw[accent1, fill=white] (-.95,3.4) circle[radius=2pt];
    \node[accent1, below] at (-.95,3.35) {$x_n$};
    \draw[->, black!50] (1.6,2.8)--(1.6,1.3) node[midway, right, text=black] {$f$};
    \draw[accent1, line width=1.2pt] (-1.4,0)--(1.4,0);
    \node[accent1, right] at (1.4,0) {$E$};
    \draw[line width=.8pt] (.3,-1)--(.3,1);
    \node[above] at (.3,1) {$F_p$};
    \node[right, font=\scriptsize] at (.3,.55) {$\times2$};
    \draw[accent1, fill=white] (-.95,0) circle[radius=2pt];
    \node at (0,-1.6) {(b)};
  \end{scope}
\end{tikzpicture}

%%FIG maps with x_n on E for beta_2 at the second step (Research_Note_2026_10_29-5.svg)
\begin{tikzpicture}[line cap=round, line join=round]
  %% (a) E with three distinct fibers: stays
  \begin{scope}
    \draw[accent1, line width=1.2pt] (-1.3,3.4)--(1.3,3.4);
    \node[accent1, above] at (0,3.4) {$Z\to E$};
    \foreach \x/\l in {-.4/{F_1}, .3/{F_2}, 1/{F_3}}{
      \draw[line width=.8pt] (\x,3.4)--(\x,2.4);
      \fill (\x,3.4) circle[radius=1.3pt];
      \node[right, font=\scriptsize] at (\x,3.15) {$1$};
      \node[below] at (\x,2.4) {$\l$};
    }
    \draw[accent1, fill=white] (-1,3.4) circle[radius=2pt];
    \node[accent1, below] at (-1,3.35) {$x_n$};
    \draw[->, black!50] (1.6,2.8)--(1.6,1.3) node[midway, right, text=black] {$f$};
    \draw[accent1, line width=1.2pt] (-1.4,0)--(1.4,0);
    \node[accent1, right] at (1.4,0) {$E$};
    \foreach \x/\l in {-.4/{F_1}, .3/{F_2}, 1/{F_3}}{
      \draw[line width=.8pt] (\x,-1)--(\x,1);
      \node[above] at (\x,1) {$\l$};
    }
    \draw[accent1, fill=white] (-1,0) circle[radius=2pt];
    \node at (0,-1.6) {(a)};
  \end{scope}
  %% (b) E with a fiber and a double cover of another fiber: comb
  \begin{scope}[shift={(4.2,0)}]
    \draw[accent1, line width=1.2pt] (-1.3,3.4)--(1.3,3.4);
    \node[accent1, above] at (0,3.4) {$Z\to E$};
    \node[right, font=\scriptsize] at (.7,2.8) {$2{:}1$};
    \foreach \x/\l in {-.3/{F_p}, .7/{F_{p'}}}{
      \draw[line width=.8pt] (\x,3.4)--(\x,2.4);
      \fill (\x,3.4) circle[radius=1.3pt];
      \node[right, font=\scriptsize] at (\x,3.15) {$1$};
      \node[below] at (\x,2.4) {$\l$};
    }
    \draw[accent1, fill=white] (-1,3.4) circle[radius=2pt];
    \node[accent1, below] at (-1,3.35) {$x_n$};
    \draw[->, black!50] (1.6,2.8)--(1.6,1.3) node[midway, right, text=black] {$f$};
    \draw[accent1, line width=1.2pt] (-1.4,0)--(1.4,0);
    \node[accent1, right] at (1.4,0) {$E$};
    \draw[line width=.8pt] (-.3,-1)--(-.3,1);
    \node[above] at (-.3,1) {$F_p$};
    \draw[line width=.8pt] (.7,-1)--(.7,1);
    \node[above] at (.7,1) {$F_{p'}$};
    \node[right, font=\scriptsize] at (.7,.55) {$\times2$};
    \draw[accent1, fill=white] (-1,0) circle[radius=2pt];
    \node at (0,-1.6) {(b)};
  \end{scope}
  %% (c) center contracted to p with x_n and x_i, a beta_2 branch transversal at p: comb
  \begin{scope}[shift={(8.4,0)}]
    \draw[accent1, line width=1.2pt] (-1.3,3.4)--(1.3,3.4);
    \node[accent1, above] at (0,3.4) {$Z\mapsto p$};
    \draw[line width=.8pt] (.5,3.4)--(.5,2.4);
    \fill (.5,3.4) circle[radius=1.3pt];
    \node[right, font=\scriptsize] at (.5,3.15) {$1$};
    \node[below] at (.5,2.4) {$\beta_2$};
    \foreach \x/\l in {-1/{x_n}, -.4/{x_i}}{
      \draw[accent1, fill=white] (\x,3.4) circle[radius=2pt];
      \node[accent1, below] at (\x,3.35) {$\l$};
    }
    \draw[->, black!50] (1.6,2.8)--(1.6,1.3) node[midway, right, text=black] {$f$};
    \draw[line width=1pt] (-1.4,0)--(1.4,0);
    \node[right] at (1.4,0) {$E$};
    \draw[line width=.8pt] plot[domain=-1.4:1.4, samples=60] (\x,{.4*(\x+.3)*(\x-.9)});
    \node[left] at (-1.4,.75) {$\beta_2$};
    \fill[accent1] (-.3,0) circle[radius=2pt];
    \node[accent1, below] at (-.3,-.1) {$p$};
    \node at (0,-1.6) {(c)};
  \end{scope}
\end{tikzpicture}

%%FIG maps with x_n on a component contracted to p for beta_3 at the third step (Research_Note_2026_10_29-6.svg)
\begin{tikzpicture}[line cap=round, line join=round]
  %% (a) beta_2 tangent to E at p and F_p: stays
  \begin{scope}
    \draw[accent1, line width=1.2pt] (-1.3,3.4)--(1.3,3.4);
    \node[accent1, above] at (0,3.4) {$Z\mapsto p$};
    \foreach \x/\l/\c in {-.2/{\beta_2}/2, .7/{F_p}/1}{
      \draw[line width=.8pt] (\x,3.4)--(\x,2.4);
      \fill (\x,3.4) circle[radius=1.3pt];
      \node[right, font=\scriptsize] at (\x,3.15) {$\c$};
      \node[below] at (\x,2.4) {$\l$};
    }
    \draw[accent1, fill=white] (-.95,3.4) circle[radius=2pt];
    \node[accent1, below] at (-.95,3.35) {$x_n$};
    \draw[->, black!50] (1.6,2.8)--(1.6,1.3) node[midway, right, text=black] {$f$};
    \draw[line width=1pt] (-1.4,0)--(1.4,0);
    \node[right] at (1.4,0) {$E$};
    \draw[line width=.8pt] plot[domain=-1.4:1.4, samples=60] (\x,{.45*\x*\x});
    \node[left] at (-1.4,.88) {$\beta_2$};
    \draw[line width=.8pt] (0,-1)--(0,1);
    \node[above] at (0,1) {$F_p$};
    \fill[accent1] (0,0) circle[radius=2pt];
    \node[accent1, below right] at (0,-.05) {$p$};
    \node at (0,-1.6) {(a)};
  \end{scope}
  %% (b) beta_2 transversal at p and F_p: comb
  \begin{scope}[shift={(4.2,0)}]
    \draw[accent1, line width=1.2pt] (-1.3,3.4)--(1.3,3.4);
    \node[accent1, above] at (0,3.4) {$Z\mapsto p$};
    \foreach \x/\l in {-.2/{\beta_2}, .7/{F_p}}{
      \draw[line width=.8pt] (\x,3.4)--(\x,2.4);
      \fill (\x,3.4) circle[radius=1.3pt];
      \node[right, font=\scriptsize] at (\x,3.15) {$1$};
      \node[below] at (\x,2.4) {$\l$};
    }
    \draw[accent1, fill=white] (-.95,3.4) circle[radius=2pt];
    \node[accent1, below] at (-.95,3.35) {$x_n$};
    \draw[->, black!50] (1.6,2.8)--(1.6,1.3) node[midway, right, text=black] {$f$};
    \draw[line width=1pt] (-1.4,0)--(1.4,0);
    \node[right] at (1.4,0) {$E$};
    \draw[line width=.8pt] plot[domain=-1.4:1.4, samples=60] (\x,{.4*(\x+.3)*(\x-.9)});
    \node[left] at (-1.4,.75) {$\beta_2$};
    \draw[line width=.8pt] (-.3,-1)--(-.3,1);
    \node[above] at (-.3,1) {$F_p$};
    \fill[accent1] (-.3,0) circle[radius=2pt];
    \node[accent1, below left] at (-.3,-.05) {$p$};
    \node at (0,-1.6) {(b)};
  \end{scope}
\end{tikzpicture}

\end{document}
